% GENERATED FILE. Edit canonical lessons, then run npm run generate:books.
% canonical-source-hash: fnv1a64:cacb7c3a5a0fe022
% canonical-lessons: TA-C268-urcakam, TA-C268-kulappam, TA-C268-punnakai, TA-C268-kanivana, TA-C268-nermaiyana

\chapter{Enthusiasm, Confusion, Smile, Gentle, Honest}
\label{ch:ta-enthusiasm-confusion-smile-gentle-honest}
\hlchaptermodality{268}

\begin{chapteropening}
\textbf{By the end of this chapter:} \emph{I can say enthusiasm, confusion, a smile, gentle and honest.}

Complete the last lesson of chapter 268: I can say enthusiasm, confusion, a smile, gentle and honest.
\end{chapteropening}

\section[\texorpdfstring{\ta{உற்சாகம்} (uṟcākam)}{உற்சாகம் (uṟcākam)}]{\ta{உற்சாகம்} (uṟcākam) — enthusiasm, excitement}

\label{lesson:TA-C268-urcakam}



\begin{quote}
Before the new one: say the Tamil for crowded, then the Tamil for easy.
\end{quote}

\subsection*{You'll want to know: \ta{உற்சாகம்}}
\textbf{\ta{உற்சாகம்}} — \emph{uṟcākam} — \textquotedblleft{}enthusiasm, excitement\textquotedblright{}.

\begin{grammarlens}[title={What you've built}]
One more word for this chapter.
\end{grammarlens}

\subsection*{Your turn}
\begin{itemize}
\raggedright
  \item \emph{Say it:} \emph{uṟcākam}
  \item \emph{Say it:} \emph{uṟcākam}, once more
  \item \emph{Say it:} say \emph{uṟcākam}
  \item \emph{Recall:} say the Tamil for crowded, then the Tamil for easy, then say \emph{uṟcākam} again
\end{itemize}

\begin{tcolorbox}[breakable,colback=teal!4,colframe=teal!35!black,title={Before you move on}]
What does \ta{உற்சாகம்} mean? (\textquotedblleft{}Enthusiasm, excitement\textquotedblright{}.) Say it once more.
\end{tcolorbox}
\section[\texorpdfstring{\ta{குழப்பம்} (kuḻappam)}{குழப்பம் (kuḻappam)}]{\ta{குழப்பம்} (kuḻappam) — confusion}

\label{lesson:TA-C268-kulappam}



\begin{quote}
Before the new one: say the Tamil for easy, then the Tamil for enthusiasm.
\end{quote}

\subsection*{You'll want to know: \ta{குழப்பம்}}
\textbf{\ta{குழப்பம்}} — \emph{kuḻappam} — \textquotedblleft{}confusion\textquotedblright{}.

\begin{grammarlens}[title={What you've built}]
One more word for this chapter.
\end{grammarlens}

\subsection*{Your turn}
\begin{itemize}
\raggedright
  \item \emph{Say it:} \emph{kuḻappam}
  \item \emph{Say it:} \emph{kuḻappam}, once more
  \item \emph{Say it:} say \emph{kuḻappam}
  \item \emph{Recall:} say the Tamil for easy, then the Tamil for enthusiasm, then say \emph{kuḻappam} again
\end{itemize}

\begin{tcolorbox}[breakable,colback=teal!4,colframe=teal!35!black,title={Before you move on}]
What does \ta{குழப்பம்} mean? (\textquotedblleft{}Confusion\textquotedblright{}.) Say it once more.
\end{tcolorbox}
\section[\texorpdfstring{\ta{புன்னகை} (puṉṉakai)}{புன்னகை (puṉṉakai)}]{\ta{புன்னகை} (puṉṉakai) — a smile}

\label{lesson:TA-C268-punnakai}



\begin{quote}
Before the new one: say the Tamil for enthusiasm, then the Tamil for confusion.
\end{quote}

\subsection*{You'll want to know: \ta{புன்னகை}}
\textbf{\ta{புன்னகை}} — \emph{puṉṉakai} — \textquotedblleft{}a smile\textquotedblright{}.

\begin{grammarlens}[title={What you've built}]
One more word for this chapter.
\end{grammarlens}

\subsection*{Your turn}
\begin{itemize}
\raggedright
  \item \emph{Say it:} \emph{puṉṉakai}
  \item \emph{Say it:} \emph{puṉṉakai}, once more
  \item \emph{Say it:} say \emph{puṉṉakai}
  \item \emph{Recall:} say the Tamil for enthusiasm, then the Tamil for confusion, then say \emph{puṉṉakai} again
\end{itemize}

\begin{tcolorbox}[breakable,colback=teal!4,colframe=teal!35!black,title={Before you move on}]
What does \ta{புன்னகை} mean? (\textquotedblleft{}A smile\textquotedblright{}.) Say it once more.
\end{tcolorbox}
\section[\texorpdfstring{\ta{கனிவான} (kaṉivāṉa)}{கனிவான (kaṉivāṉa)}]{\ta{கனிவான} (kaṉivāṉa) — gentle, kind}

\label{lesson:TA-C268-kanivana}



\begin{quote}
Before the new one: say the Tamil for confusion, then the Tamil for a smile.
\end{quote}

\subsection*{You'll want to know: \ta{கனிவான}}
\textbf{\ta{கனிவான}} — \emph{kaṉivāṉa} — \textquotedblleft{}gentle, kind\textquotedblright{}.

\begin{grammarlens}[title={What you've built}]
One more word for this chapter.
\end{grammarlens}

\subsection*{Your turn}
\begin{itemize}
\raggedright
  \item \emph{Say it:} \emph{kaṉivāṉa}
  \item \emph{Say it:} \emph{kaṉivāṉa}, once more
  \item \emph{Say it:} say \emph{kaṉivāṉa}
  \item \emph{Recall:} say the Tamil for confusion, then the Tamil for a smile, then say \emph{kaṉivāṉa} again
\end{itemize}

\begin{tcolorbox}[breakable,colback=teal!4,colframe=teal!35!black,title={Before you move on}]
What does \ta{கனிவான} mean? (\textquotedblleft{}Gentle, kind\textquotedblright{}.) Say it once more.
\end{tcolorbox}
\section[\texorpdfstring{\ta{நேர்மையான} (nērmaiyāṉa)}{நேர்மையான (nērmaiyāṉa)}]{\ta{நேர்மையான} (nērmaiyāṉa) — honest}

\label{lesson:TA-C268-nermaiyana}



\begin{quote}
Before the new one: say the Tamil for a smile, then the Tamil for gentle.
\end{quote}

\subsection*{You'll want to know: \ta{நேர்மையான}}
\textbf{\ta{நேர்மையான}} — \emph{nērmaiyāṉa} — \textquotedblleft{}honest\textquotedblright{}.

\begin{grammarlens}[title={What you've built}]
One more word for this chapter.
\end{grammarlens}

\subsection*{Your turn}
\begin{itemize}
\raggedright
  \item \emph{Say it:} \emph{nērmaiyāṉa}
  \item \emph{Say it:} \emph{nērmaiyāṉa}, once more
  \item \emph{Say it:} say \emph{nērmaiyāṉa}
  \item \emph{Recall:} say the Tamil for a smile, then the Tamil for gentle, then say \emph{nērmaiyāṉa} again
\end{itemize}

\begin{tcolorbox}[breakable,colback=teal!4,colframe=teal!35!black,title={Before you move on}]
What does \ta{நேர்மையான} mean? (\textquotedblleft{}Honest\textquotedblright{}.) Say it once more.
\end{tcolorbox}
